% ===== 原文 pp.19–21 =====
% tr-03 块，第 2 片：
%   §3.2.2 收尾（含图 9）+ §3.2.3 四穿刺球面 + §3.2.4 一般性讨论
%   图 9（p.19）、图 10（p.20）；tag 3.7–3.11
% 本片开头一句是原文 p.18 的最后一行（跨页到 p.19），为不切断句子整体
% 移到此处，与 tr03a.tex 末句直接相接。

% [原文 p.18 末句，跨页至 p.19]
我们可以利用 $\SU(2)$ 味对称性给出质量项，同时仍保持 $\cN=2$ 超对称性。所得理论通常称为 $\cN=2^*$ 理论。

% [原文 p.19]

\begin{center}
\begin{minipage}{\textwidth}
\centering
\includegraphics[width=0.45\textwidth]{fig9.png}\par\vspace{2pt}
\figtitle{9}{带穿刺点的环面可由三穿刺球面粘合两个点得到。}
\end{minipage}
\end{center}

注意，我们得到的这个规范理论的单圈 $\beta$ 函数为零。如此之高的超对称性保证 $\beta$ 函数即使在非微扰层面也为零。因此，复化耦合 $\tau$ 在该规范理论中始终是一个无量纲参数。

% [Tutor Note: p.19，“零单圈 $\beta$ 函数”可以逐条验一下。用 $\cN=2$ 的单圈
%  系数 $b_0=2N_c-N_f$（$N_f$ 按基础超多重态计），本文只用 $N_c=2$：
%  单穿刺环面情形，$Q_{aiu}$ 对固定的 $u$ 在 $(a,i)$ 上有 4 个分量，把
%  两个指标认作同一个规范指标后 $2\otimes2=3\oplus1$，于是规范荷流下的
%  物质是伴随手征场 $A_{Iu}$（$I=1,2,3$）——连同 $u=1,2$ 正好是一个
%  伴随超多重态，把伴随折合成 $N_f=2N_c=4$，得 $b_0=2\cdot2-4=0$；
%  四穿刺球面情形，$Q_{aiu}$ 与 $Q'_{usx}$ 各在（$a,i$ 或 $s,x$）上给出
%  4 个 $u$ 二重态，即各 2 个超多重态，两个加起来 $N_f=4$，同样
%  $b_0=4-4=0$。这正是 “任何一条圆柱面对应的规范群 $\beta$ 函数都为零” 的
%  来源，也是 §3.2.4 能说 “所有耦合常数都是真正的无量纲参数” 的原因。]

$\cN=4$ 超杨–米尔斯理论有一个重要性质，即它的 S-对偶：规范群为 $G$、耦合常数为 $\tau$ 的理论，与规范群为对偶群 $G^\vee$、耦合常数为 $-1/\tau$ 的理论等价。我们只用到 $G=\SU(2)$ 的情形；这时对偶群 $G^\vee$ 恰好与原群相同，于是 $G=G^\vee=\SU(2)$。

% [Tutor Note: p.19，S-对偶 $G\leftrightarrow G^\vee$、$\tau\to-1/\tau$。
%  $G^\vee$ 指 Langlands（Montonen–Olive）对偶群。原文说 $\SU(2)$ 的对偶群
%  “恰好与原群相同”，是在李代数/局域算符层面说的：严格地讲
%  $\SU(2)$ 的对偶群是 $\SU(2)/\Z_2=\SO(3)$，两者李代数同为
%  $\mathfrak{su}(2)$，差别只体现在允许的线算符（Wilson 圈与 't Hooft 圈）
%  谱上，不影响本节关心的拉格朗日量数据与 $\tau$ 的变换。另外注意
%  $\tau\to-1/\tau$ 只是 $S$ 变换；$T$ 变换 $\tau\to\tau+1$ 在 (3.6) 那里
%  已作为同一圆柱面的大规范变换处理掉了（$\theta\sim\theta+2\pi$）。]

在我们的规范理论几何构造中，这一点有很好的解释：从一个三穿刺球面加一个圆柱面出发，我们可以用两种方式构造出模参数为 $\tau$ 的单穿刺环面，即分别用参数为 $\tau$ 的圆柱面和参数为 $-1/\tau$ 的圆柱面。

相应地，对这个单穿刺环面，按我们的规则可以配上两个规范理论：规范群为 $\SU(2)$、复化规范耦合为 $\tau$ 的 $\cN=4$ 理论，或同一理论但规范耦合为 $-1/\tau$。S-对偶保证这两个理论是同一个理论。

\subsubsection{四穿刺球面}

接下来取两个三穿刺球面，用一个参数为 $\tau$ 的圆柱面把一个球面上的一个穿刺点与另一个球面上的一个穿刺点连接起来。这样就得到一个四穿刺球面。

用规范理论的语言说，我们有两个三基表示 $Q_{aiu}$ 与 $Q'_{usx}$。指标 $a,i,s,x=1,2$ 对应四个 $\SU(2)$ 对称性，而有一个动力学 $\SU(2)$ 规范多重态作用在指标 $u=1,2$ 上。单圈 $\beta$ 函数为零。同样，$\beta$ 函数即使在非微扰层面也为零，复化耦合常数是该理论真正的无量纲参数。

我们可以把手征物质场重新组织为 $q_{uI}$，其中 $u=1,2$、$I=1,\dots,8$。超势耦合为
\begin{equation}
  \delta^{IJ} q_{uI} \Phi_{uv} q_{vJ}
  \tag{3.7}
\end{equation}
其中 $\Phi$ 是 $\SU(2)$ 规范多重态中的伴随手征标量。指标 $uv$ 是对称的，因此 $\delta_{IJ}$ 也是对称的。这意味着味对称性是 $\SO(8)$。这个理论通常称为带 $N_f=4$ 个味的 $\SU(2)$ 规范理论。这四个 $\SU(2)$ 是 $\SO(8)$ 的%
% [原文 p.20]
子群：
\begin{equation}
  \SU(2)_a \times \SU(2)_i \times \SU(2)_s \times \SU(2)_x \subset \SO(8)
  \tag{3.8}
\end{equation}
而手征物质场则按
\begin{equation}
  2_a \otimes 2_i \oplus 2_s \otimes 2_x = 8_V
  \tag{3.9}
\end{equation}
变换，其中 $8_V$ 是 $\SO(8)$ 味对称性的矢量表示。

% [Tutor Note: p.19–p.20，从 (3.7) 到 (3.9)。(3.7) 里 $u,v$ 走
%  $q_{uI}\Phi_{uv}q_{vJ}$，$\Phi_{uv}$ 对 $uv$ 对称（它正是 $\SU(2)$ 伴随
%  标量，作为 $2\otimes2$ 的对称部分即三重态），所以把 $I,J$ 缩并起来的
%  $\delta^{IJ}$ 也必须对称。这一点有可检验的后果：八个 $I$ 之间带一个
%  对称的不变双线性型，说明这个 8 维表示是正交型的，于是味对称群是
%  $\SO(8)$（辛群 $\USp(8)$ 的不变形式是反对称的，故不可能是它）。
%  (3.9) 的维数核对：$2_a\otimes2_i$ 与 $2_s\otimes2_x$ 各 4 维，
%  $4+4=8$，正是 $\SO(8)$ 的矢量表示 $8_V$ 限制到 $\SU(2)^4$ 的分解。]

\begin{center}
\begin{minipage}{\textwidth}
\centering
\includegraphics[width=0.55\textwidth]{fig10.png}\par\vspace{2pt}
\figtitle{10}{四穿刺球面可以有三种不同的分解方式。}
\end{minipage}
\end{center}

现在，四穿刺球面可以由两个三穿刺球面以三种不同方式得到，见图 10。我们已经描述过其中一种。另外两种同样给出带 $N_f=4$ 个味的 $\SU(2)$ 理论，但 $\SO(8)$ 表示不同、耦合常数也不同。具体地说，物质场处于
\begin{equation}
  2_a \otimes 2_s \oplus 2_i \otimes 2_x = 8_S
  \tag{3.10}
\end{equation}
耦合为 $q' = 1/q$；而物质场处于
\begin{equation}
  2_a \otimes 2_x \oplus 2_s \otimes 2_i = 8_C
  \tag{3.11}
\end{equation}
耦合为 $q'' = 1-q$。这里 $8_S$ 与 $8_C$ 分别是 $\SO(8)$ 的正、负手征旋量，而我们用了指数化的复化耦合常数 $q=\ee^{2\pi\ii\tau}$、$q'=\ee^{2\pi\ii\tau'}$、$q''=\ee^{2\pi\ii\tau''}$。同样，保证这三种描述等价的 S-对偶早已为人所知。

% [Tutor Note: p.20，(3.9)–(3.11) 与图 10。(3.9)(3.10)(3.11) 只是把
%  四个指标 $a,i,s,x$ 的三种配对方式写出来：(3.9) 配 $(ai)(sx)$、
%  (3.10) 配 $(as)(ix)$、(3.11) 配 $(ax)(si)$，与图 10 的三幅图一一对应；
%  而 $8_V/8_S/8_C$ 的互换正是 $\SO(8)$ 的外部自同构（triality），这解释了
%  为什么三种描述给出同一个理论、却把同一批物质场写成不同的 $\SO(8)$
%  表示。另外 $q'=1/q$、$q''=1-q$ 要按四穿刺球面的交比来读：球面上
%  四个点的三种配对对应的交比正是 $z$、$1/z$、$1-z$（anharmonic 群
%  的三个元素）。若把 $q'=1/q$ 与 $q=\ee^{2\pi\ii\tau}$ 一起读成
%  $\ee^{2\pi\ii\tau'}=\ee^{-2\pi\ii\tau}$，会得到 $\tau'=-\tau$，这与
%  p.19 的 S-对偶 $\tau\to-1/\tau$ 并不是同一个变换——两种说法属于不同
%  的“耦合”定义（一个按曲面交比、一个按规范耦合），引用具体公式时以
%  $\tau$ 上的 $SL(2,\Z)$/S-对偶变换为准即可。]

\subsubsection{一般性讨论}

让我们把迄今引入的内容复述一遍。第一，我们给出了构造一个 $\cN=2$ 规范理论的方法。数据被编码为三穿刺球面，以及连接各对穿刺点的圆柱面。容易验证，与任何一条圆柱面对应的每个规范群，其单圈 $\beta$ 函数都为零。$\cN=2$ 超对称随即保证所有 $\beta$ 函数即使在非微扰层面也为零，因此所有规范耦合常数都保持为该理论真正的无量纲参数。那些没有通过圆柱面连接起来的穿刺点提供 $\SU(2)$ 味对称性。有 $k$ 个穿刺点，就有 $\SU(2)^k$ 味对称性；如有需要，我们还可以打开相应的 $k$ 个质量参数。至此，对于给定的带穿刺点的 Riemann 曲面，在指定了到三穿刺球面的分解之后，我们就得到了一个拉格朗日量场论。

% [Tutor Note: p.20，“$k$ 个穿刺点 $\Rightarrow$ $\SU(2)^k$ 味对称性，
%  可打开 $k$ 个质量参数”。数目对得上是因为由 p.17 的讨论质量矩阵必须
%  满足 $[m,m^\dagger]=0$（即 $m$ 可对角化），能开的只是味对称群
%  Cartan 方向上的参数；单个 $\SU(2)$ 的 Cartan 只有一维，故每个穿刺点
%  只贡献一个（复）质量参数，$k$ 个穿刺点合起来是 $k$ 个，而不是 $3k$ 个。]

第二，我们考虑过单穿刺环面；它可以由两个三穿刺球面以两种方式构造出来。相应的两个规范理论都是规范群为 $\SU(2)$、带一个解耦超多重态的 $\cN=4$ 超杨–米尔斯理论，只是耦合分别为 $\tau$ 与 $-1/\tau$。已知它们互为 S-对偶。因此，单穿刺环面实际上对应的是单个理论。

% [Tutor Note: p.20，“a torus with one puncture can be constructed from
%  TWO three-punctured spheres in two ways”。这句与 p.19 的陈述冲突：
%  p.19 明确说单穿刺环面是 “从一个三穿刺球面加一个圆柱面、以两种方式”
%  构造的。用亏格/穿刺点计数也能判明前者不可能：两个三穿刺球面粘一对
%  穿刺点，得 3+3-2=4 个穿刺点、亏格 0 的球面，不会得到亏格 1、一个穿刺点的
%  环面。已按 700 dpi 复核原文确实印作 “two”，故应为 “a three-punctured
%  sphere and a cylinder”（“two ways” 指圆柱面参数取 $\tau$ 或 $-1/\tau$，
%  与 p.19 一致）。同一句另有脱字，规范说法应是 “we considered how a torus
%  ... can be constructed”。正文照原文排。]

第三，我们考虑过四穿刺球面。它可以由两个三穿刺球面以三种方式构造出来。相应的三个规范理论都是带 $N_f=4$ 个味的 $\SU(2)$ 理论，但超多重态处于不同的 $\SO(8)$ 味表示、耦合常数也各不相同。同样，已知它们互为 S-对偶。因此，四穿刺球面实际上对应的是单个理论。

% [原文 p.21]

一般而言，给定一个带若干穿刺点的单个 Riemann 曲面，把它切成三穿刺球面与圆柱面的方式有多种。每一种这样的分解都给出一个不同的 $\cN=2$ 超对称规范理论拉格朗日量。然而，把上面回顾的两条 S-对偶结合起来，任何两种这样的分解都能相互联系起来。因此我们得出结论：带若干穿刺点的单个 Riemann 曲面，实际上对应的是单个量子场论，各种不同的拉格朗日量只是它的化身（avatar）而已。

% [Tutor Note: p.21，末句。“两条 S-对偶”指：单穿刺环面给的
%  $\tau\to-1/\tau$（p.19），以及四穿刺球面给的三个描述之间的对偶
%  （p.20）。论证的结构是：每种分解给一个拉格朗日量，任意两种分解之间只差
%  有限次 S-对偶变换，于是它们互为 S-对偶描述；S-对偶是理论层面的等价，
%  不是拉格朗日量层面的等价，所以 “不同的拉格朗日量只是同一个理论的化身”。
%  这句是全节的目标：既然理论只依赖曲面本身，就可以把配分函数看作
%  曲面（连同穿刺点与分解数据）的函数，这正是后面 §5 把 4d 配分函数
%  与 2d Liouville 理论对应起来的出发点。]

% 边界提醒（非批注）：本文件末句 “…are just avatars.” 之后紧接的 p.21
% “4 Two supersymmetric backgrounds” 标题属 tr-04，本块不得写入。
